% 原书分章 PDF 第 3—4 页：第 28.2 节导言与第 28.2.1 节。
\section{TMLE}

如本书其他部分所述，TMLE 为目标参数生成渐近线性的代入估计量。TMLE 对目标参数的初始估计作扰动，从而得到能够正确权衡偏差与方差的估计。即使似然的各组成部分采用这类数据自适应方法估计，TMLE 仍然渐近线性，并具有已知的影响曲线。

\subsection{特定处理下均值的 TMLE}

$\Psi$ 在 $P$ 处的有效影响曲线为
\[
 D^*(P)(O)=\frac{A}{\bar g(W)}(Y-\bar Q(W))
 +\bar Q(W)-\Psi(P)
\]
（van der Laan 和 Robins，2003）。注意，$\Psi(P)$ 仅通过 $\bar Q(P)$ 及 $W$ 的概率分布 $Q_W(P)$ 依赖 $P$。令 $Q(P)=(Q_W(P),\bar Q(P))$，并令
$Q(\mathcal M)=\{Q(P):P\in\mathcal M\}$ 为其模型空间。我们也将目标参数记为 $\Psi:Q(\mathcal M)\to\mathbb R$，表示将参数空间 $Q(\mathcal M)$ 中的一个 $Q$ 映射为数值的映射；这里略微滥用了记号，对这一映射仍使用相同的 $\Psi$。类似地，也将 $D^*(P)$ 记为 $D^*(Q,G)$。有效影响曲线 $D^*(P)$ 满足展开
$\Psi(P)-\Psi(P_0)=-P_0D^*(P)+R_\psi(P,P_0)$，其中
\[
 R_\psi(P,P_0)=P_0\frac{\bar g-\bar g_0}{\bar g}
 (\bar Q-\bar Q_0).
\]

令 $\psi_n^*=\Psi(Q_n^*)$ 为 $\psi_0$ 的一个 TMLE，使其在 $P_0$ 处渐近线性，影响曲线为 $D^*(P_0)$。这个 TMLE 可按如下方式定义。令 $\bar Q_n^0$ 为 $\bar Q_0$ 的初始估计量，$\bar g_n$ 为 $g_0$ 的估计量，
\[
 L(\bar Q)(O)=-I(A=1)(Y\log\bar Q(W)
 +(1-Y)\log(1-\bar Q(W))
\]
为 $\bar Q_0$ 的对数似然损失函数。取子模型
\[
 \mathrm{Logit}\,\bar Q_n^0(\epsilon))
 =\mathrm{Logit}\,\bar Q_n^0+\epsilon H(\bar g_n),
 \qquad H(\bar g_n)=A/\bar g_n(W).
\]
令 $\bar Q_n^1=\bar Q_n^0(\epsilon_n^0)$，其中
$\epsilon_n^0=\argmin_\epsilon P_nL(\bar Q_n^0(\epsilon))$；并令
$\psi_n^*=\Psi(Q_n^1)$，其中 $Q_n^1=(\bar Q_n^1,Q_{W,n})$，
$Q_{W,n}$ 为 $W_1,\ldots,W_n$ 的经验分布。令 $P_n^*$ 为与 $Q_n^*$ 相容的概率分布。
